\section{What the results say about nonsmooth BFGS}\label{art:geometry}

The mixed ridge $F(u,v)=u^2+|v|$ isolates the difficulty left by a
subsequence approaching the nonsmooth line. Reaching $v=0$ does not eliminate
the horizontal error $u^2$. Yet the vertical gradient jump does not remove
the direction that can reduce this error. Along $v=0$ the limiting gradients are $(2u,1)$ and
$(2u,-1)$. Their difference is vertical. At a point with $u\ne0$, the
horizontal vector $(-u,0)$ has the same strictly negative pairing $-2u^2$
with both gradients. The inward radial model also has a small error:
\begin{equation}\label{art:mixed-inward-remainder}
 F((1-\tau)(u,v))-F(u,v)+\tau\nabla F(u,v)^T(u,v)
       =\tau^2u^2\qquad(v\ne0,\ 0\le\tau\le1).
\end{equation}
Thus the nonsmooth part is exactly linear under inward scaling, while the
smooth part contributes a quadratic remainder. Section~\ref{art:global}
uses this property to prove convergence of the whole sequence. The inward
segment is a comparison used in the proof; it is not a direction BFGS is
assumed to take. No eventual eigenvector or eigenvalue description is needed.

\begin{figure}[htbp]
\centering
% Vector construction from explicit gradients; no trajectory inference.
\begin{tikzpicture}[x=1.7cm,y=1.5cm,>=Stealth,font=\small]
\begin{scope}[shift={(-1.9,0)}]
\draw[->,gray] (-.8,0)--(1.55,0);\draw[->,gray] (0,-1.35)--(0,1.35);
\draw[line width=1.1pt] (1,-1)--(1,1);
\draw[->] (0,0)--(1,1);\draw[->] (0,0)--(1,-1);
\fill (1,1) circle (1.4pt) node[right] {$(1,1)$};
\fill (1,-1) circle (1.4pt) node[right] {$(1,-1)$};
\draw[->,thick,dashed] (0,0)--(-.65,0) node[below] {descent};
\node at (.4,-1.65) {(a) $u^2+|v|$, $u=1/2$};
\node[above left] at (0,0) {$0$};
\end{scope}
\begin{scope}[shift={(1.8,0)}]
\draw[->,gray] (-1.35,0)--(1.4,0);\draw[->,gray] (0,-1.35)--(0,.55);
\fill[gray!15] (1,0)--(-1,-.125)--(.3184814276,-1.0252321197)--cycle;
\draw[thick] (1,0)--(-1,-.125)--(.3184814276,-1.0252321197)--cycle;
\draw[->] (0,0)--(.3184814276,-1.0252321197);
\fill (1,0) circle (1.4pt) node[above right] {$g_0$};
\fill (-1,-.125) circle (1.4pt) node[below left] {$g_1$};
\fill (.3184814276,-1.0252321197) circle (1.4pt) node[right] {$g_2$};
\fill[white] (0,0) circle (1.7pt);\draw (0,0) circle (1.7pt);
\node[above] at (0,.04) {$0$};
\node at (0,-1.65) {(b) Certified nonstationary seed};
\end{scope}
\end{tikzpicture}

\caption{Gradient space. Left: the mixed ridge at $u=1/2$ has a vertical
jump and a common horizontal descent direction. Right: the particular
three-gradient seed used in Section~\ref{art:failure}. Its convex hull
excludes zero, but its affine hull is the plane. The vertices on the right
are rendered from the certified root's rational center; the mathematical
definition uses the isolated root. Arrows indicate gradients, not BFGS steps.}
\label{art:gradient-geometry}
\end{figure}

The right panel asks a different question: can a genuine full-BFGS run
contract around a point where descent still exists? The bad corner is specific. Its slopes
are $g_0=(1,0)$, $g_1=(-1,-1/8)$ and $g_2=(c_\circ,d_\circ)$, where $(c_\circ,d_\circ)$ is
defined exactly in Appendix~\ref{art:seed_certificate}. They are affinely
independent, so no nonzero direction pairs equally with all three. Yet
$(-1,16)$ has negative pairing with each, and zero is outside their convex
hull. The corner is therefore nonstationary. Section~\ref{art:failure}
constructs a three-step cycle and verifies that the stated line search
produces it, including all rejected trials. This additional dynamical proof
is what makes the corner a counterexample.

This distinction between affine and convex balance helps interpret the
examples. For a nonempty finite family of slopes $a_i$, equality $a_i^Tv=\kappa<0$ for every $i$ asks
for a common rate of descent, whereas $a_i^Tv<0$ asks only for descent.
The former fails when zero belongs to their affine hull; the latter fails
when zero belongs to their convex hull. Indeed, a linear functional constant
on the slopes is constant on their affine hull. Conversely, if zero is
outside that affine hull, its nearest point gives a separating functional
constant on the hull; changing its sign gives the negative common value.
The convex version follows from the same closest-point argument for the
convex hull. These are elementary separation facts, not a classification
of BFGS behavior.

The positive conclusions differ in their geometry and in the locality of the
guarantee. Table~\ref{art:positive-outcomes} distinguishes them before the
proofs. In every row the quantifier is over infinite admissible runs with
arbitrary positive-definite initialization. An admissible run uses actual
gradients and satisfies Armijo--weak-Wolfe at its accepted points; it does not
assert that an arbitrary line-search implementation always finds such points.

\begin{table}[htbp]
\centering
\begin{tabular}{@{}p{.40\textwidth}p{.22\textwidth}p{.29\textwidth}@{}}
\toprule
Objective geometry & Starting scope & Conclusion\\
\midrule
Centered smooth--sublinear, Theorem~\ref{art:centered}
& Compact initial sublevel in the structural region
& Optimal values; unique minimizer gives whole-sequence convergence\\
Coercive degree-one homogeneous, including $\ell_1$ and $\ell_\infty$ norms,
Theorem~\ref{art:homogeneous}
& Every starting point
& Whole sequence converges to zero\\
Finite polyhedral smooth-piece partition, including squared loss plus
$\lambda\norm{x}_1$ ($\lambda>0$), Theorem~\ref{art:polyhedral}
& Common low sublevel near the unique minimizer
& Whole sequence converges to the minimizer\\
Nondegenerate planar smooth maximum, Corollary~\ref{art:planar-max}
& Local basin under the stated contour hypotheses
& Whole sequence converges to the unique minimizer\\
\bottomrule
\end{tabular}
\caption{What the positive theory guarantees. The benchmark $u^2+|v|$ has
compact sublevels everywhere, so the first row gives global convergence.
The results concern accepted runs, without a presumed matrix limit.}
\label{art:positive-outcomes}
\end{table}

The first two rows make precise global assertions; they do not claim global
convergence for every affine maximum. Indeed, the coercive polyhedral
realization of the wrong corner in Corollary~\ref{art:failure_polyhedral}
has a positive limiting objective gap. General affine offsets can place the
bad corner away from the true optimizer, whereas homogeneity centers all
linear pieces at the origin.

Local behavior raises a further distinction. Finitely many smooth formulas
can reproduce bad configurations at infinitely many shrinking scales.
A finite polyhedral partition near the optimizer controls the inward model
because each remaining cell contains the optimizer after localization.
Curved interfaces require a different argument: the planar example in
Section~\ref{art:local} has a radial obstruction, yet a moving contour field
proves convergence. These are complementary sufficient mechanisms and
specific realized failures, rather than a classification by jump span.
